\documentclass[main.tex]{subfiles}
\begin{document}
	
\chapter{Works of J.-G. Jiang (I)} 
	\label{ch5}

\section{Normal Covers and $D\omega^2$-refinable Spaces}
\label{sec:5-1}

{\theorem{\label{th:5.1.1}}
For any open cover $\mathscr{U}$ of $X$, the following are equivalent.
	
\textup{(i)}  $\mathscr{U}$ is a normal cover.

\textup{(ii)} $\mathscr{U}$ has a $\sigma$-disjoint functionally open refinement $\mathscr{V} = \bigcup_{n<\omega} \mathscr{V}_n$ such that each $\bigcup \mathscr{V}_n$ is functionally open.

\textup{(iii)}  $\mathscr{U}$ has a $\sigma$-star finite functionally open refinement 
$\mathscr{V} = \bigcup_{n<\omega} \mathscr{V}_n$ such that each $\bigcup \mathscr{V}_n$ is functionally open.
}

\begin{proof}
(i) $\Rightarrow$ (ii). Suppose $\mathscr{U}$ is normal. By Theorem \ref{th:2.1.6}, $\mathscr{U}$ has a functionally open refinement $\mathscr{V} = \bigcup_{n<\omega} \mathscr{V}_n$ where each $\mathscr{V}_n$ is discrete, so is disjoint, and by Lemma \ref{lem:1.3.8}, $\bigcup \mathscr{V}_n$ is functionally open.

$(ii) \Rightarrow (iii)$. Trivial.

$(iii) \Rightarrow (i)$. Suppose $\mathscr{U}$ is an open cover of $X$ and 
$\mathscr{V} = \bigcup_{n<\omega} \mathscr{V}_n$ is a functionally open refinement 
of $\mathscr{U}$ such that each 
$\mathscr{V}_n=\{V_{n,\alpha}: \alpha\in\Delta_n\}$ is star-finite and 
$\bigcup \mathscr{V}_n$ is functionally open. By Theorem \ref{th:2.1.3}, 
$\mathscr{V}' = \{\bigcup \mathscr{V}_n: {n<\omega}\}$ has a star-finite 
functionally open refinement $\mathscr{L}$. Since $\mathscr{L}$ is locally finite, 
$\mathscr{V}'$ is a normal cover of $X$. Let $\mathscr{W}$ 
be a normal cover of \(X\) such that 
\(\overline{\mathscr{W}}\) refines \(\mathscr{V}'\). 
By Lemma \ref{lem:1.1.5}, 
there exists a locally finite functionally open cover
\(\mathscr{H}=\{H_n: n < \omega\) such that 
\(\overline{H_n}\subset\cup\mathscr{V}_n\) for each \(n<\omega\).
For each \(n<\omega\), let \(\mathscr{G}_n=\{H_n\cap V_{n,\alpha}: \alpha\in\Delta_n\}\) 
where each set \(H_n\cap V_{n,\alpha}\) is
functionally open. To see that \(\mathscr{G}_n\) is locally finite, let \(x\in X\). If
\(x\notin\cup \mathscr{V}_n\), then 
\(X\setminus\overline{H}_n\in\operatorname{top}(x)\) 
which intersects no member of \(\mathscr{G}_n\). 
If \(x\in V_{n,\beta}\) for some \(\beta\in\Delta_n\), then the set
\(\{\alpha\in\Delta_n: V_{n,\beta}\cap (H_n\cap V_{n,\alpha})\neq\emptyset\}\) 
is finite because \(\mathscr{V}_n\) is
star-finite. Hence each \(\mathscr{G}_n\) is locally finite. 
By Theorem \ref{th:2.1.3}, 
there exists a countable locally finite functionally open cover 
\(\mathscr{S}=\{S_n: n<\omega\}\)
such that \(S_n\subset\bigcup\mathscr{G}_n \forall \, n\). 
Let
$$\mathscr{R}=\{S_n\cap H_n\cap V_{n,\alpha}:  \alpha\in\Delta_n, n < \omega\}.$$

\noindent Then \(\mathscr{R}\) is a functionally open refinement of \(\mathscr{U}\). To see
that \(\mathscr{R}\) is locally finite, let \(x\in X\). Since \(\mathscr{S}\) is locally finite, 
\(\exists \, G_0\in\operatorname{top}(x)\) such that 
$\operatorname{ind}(\mathscr{S})_{G_0} = \{n_1, \cdots, n_k\}$
 is finite. For each \(i\le k\), \(\mathscr{S}_{n_i}\)
is locally finite, so \(\exists G_i\in\operatorname{top}(x)\) such that the set
$$D_i=\{\alpha\in\Delta_{n_i}: G_i\cap H_{n_i}\cap V_{n_i,\alpha}\neq\emptyset\}$$

\noindent is finite. Let \(G=\bigcap_{i=0}^k G_i\), and let
\[C=\{(\alpha, n):  \alpha\in \Delta_n, n < \omega,  
G\cap(S_{n}\cap H_n\cap V_{n,\alpha})\neq\emptyset\}.\]
\noindent It is easy to see that \(C\subset\bigcup_{i=1}^k (D_i\times\{n_i\})\). Thus \(C\) is finite,
so \(\mathscr{R}\) is locally finite. By Theorem \ref{th:2.1.6}, \(\mathscr{U}\) is a normal cover.
\end{proof}

{\theorem{\label{th:5.1.2}}
 For any open cover \(\mathscr{U}\) of \(X\) the following
are equivalent:

\textup{(i)} \(\mathscr{U}\) is a normal cover.

\textup{(ii)} There is a sequence $\{\mathscr{E}_n: n < \omega\}$ of semi-open covers of $X$ 
such that $\mathscr{E}_0$ refines $\mathscr{U}$ and each $\mathscr{E}_{n+1}$ is 
a star refinement of $\mathscr{E}_n$.

\textup{(iii)} There is a sequence $\{\mathscr{E}_n: n < \omega\}$ of semi-open covers of $X$ 
such that $\mathscr{E}_0$ refines $\mathscr{U}$ and each $\mathscr{E}_{n+1}$ is 
a point-star refinement of $\mathscr{E}_n$.
}

\begin{proof}
 (i) $\Rightarrow$ (ii) $\Rightarrow$ (iii). Clear

(iii) $\Rightarrow$ (i). Assume (iii). For each $n < \omega$, let 
$$\mathscr{U}_n = \{\operatorname{st}(x, \mathscr{E})^{\circ} : x \in X\}.$$
Then each $\mathscr{U}_n$ is an open cover of $X$. To see that $\mathscr{U}_0$ refines $\mathscr{U}$, let $x \in X$. There exists $E \in \mathscr{E}_0$ and 
$V \in \mathscr{U}$ such that $\operatorname{st}(x, \mathscr{E}_1) \subset E \subset V$. 
Then $\operatorname{st}(x, \mathscr{E}_1)^{\circ} \subset V$, 
so $\mathscr{U}_0$ refines $\mathscr{U}$.
To see that $\mathscr{U}_{n+1}$ point-star refines $\mathscr{U}_n$, let $x \in X$. 
It suffices to prove that 
$\mathrm{st}(x, \mathscr{U}_{n+1}) \subset \mathrm{st}(x, \mathscr{E}_{n+1})^{\circ}$. 
Suppose $y \in \mathrm{st}(x, \mathscr{U}_{n+1})$, then 
$x, y \in \operatorname{st}(u, \mathscr{E}_{n+2})^{\circ}$ 
for some $u \in X$. Let $E \in \mathscr{E}_{n+1}$ be such that 
$\mathrm{st}(u, \mathscr{E}_{n+2}) \subset E$. Then $x, y \in E^\circ$, so 
$y \in E^{\circ} \subset \mathrm{st}(x, \mathscr{E}_{n+1})^{\circ}$. 
This completes the proof of 
$\mathrm{st}(x, \mathscr{U}_{n+1}) \subset \mathrm{st}(x, \mathscr{E}_{n+1})^{\circ}$. 
Hence $\mathscr{U}_{n+1}$ point-star refines $\mathscr{U}_n$ and  $\mathscr{U}$ is normal.
\end{proof}



\subfile{ch5_2}

\end{document}
